\documentclass[a4paper,10pt]{article}
\usepackage[utf8]{inputenc}

% ==========================================
%      CONFIGURAZIONE UTENTE (MODIFICA QUI)
% ==========================================

% 1. Dati del Corso e Università
\newcommand{\Universita}{Università di PongiBonsi}
\newcommand{\Dipartimento}{Dipartimento di Informatica \\ Corso di Laurea Triennale in Informatica}
\newcommand{\InfoCorso}{Corso 1° anno - 12 CFU}
\newcommand{\TitoloCorso}{Analisi Matematica}
\newcommand{\Sottotitolo}{Appunti aggiornati}
\newcommand{\AnnoAccademico}{Anno Accademico 2020/2021}
\newcommand{\LogoImmagine}{marchio_unipi_pant541.png} % Assicurati che il file esista

% 2. Dati Persone (Copia e incolla dai tuoi appunti)
\newcommand{\Professori}{
    {\large\textbf{Professore:}}\\ \large{Prof. Mattia Talpo} \\
    \vspace{0.5em}
    {\large\textbf{Professore:}}\\ \large{Prof. Elisabetta Chiodaroli}
}
\newcommand{\Autori}{
    {\large\textbf{Autore:}}\\ \large{Joseph Zucchelli} \\
    \vspace{0.5em}
    
}

% 3. Colore Principale (ROSSO)
\usepackage{xcolor}
\definecolor{mainRed}{RGB}{200, 0, 0} % Rosso scuro personalizzabile
% ==========================================

% ---- Packages ---- 
\usepackage{amsmath, amsfonts, amsthm, amssymb}
\usepackage{graphicx}
\usepackage{makecell}
\usepackage{wrapfig}
\usepackage{caption}
\usepackage{subcaption}
\usepackage{geometry}
\usepackage{fancyhdr}
\usepackage{hyperref}
\usepackage[italian, english]{babel}

\graphicspath{ {./images/} }
\geometry{letterpaper,top=2cm,bottom=2cm,left=3cm,right=3cm,marginparwidth=1.75cm}

% Setup Link col colore ROSSO
\hypersetup{
    colorlinks,
    citecolor=black,
    filecolor=black,
    linkcolor=mainRed, % Link interni rossi
    urlcolor=mainRed  
}

% ---- Teoremi e Definizioni ----
\theoremstyle{remark}      \newtheorem{note}{Note}[subsection]
\theoremstyle{definition}  \newtheorem{observation}{Osservazione}[subsection]
\theoremstyle{plain}       \newtheorem{definition}{Definizione}[subsection]
\theoremstyle{definition}  \newtheorem{example}{Esempio}[subsection]
\theoremstyle{definition}  \newtheorem{demostration}{Dimostrazione}[subsection]
\theoremstyle{definition}  \newtheorem{theorem}{Teorema}[subsection]
\theoremstyle{plain}       \newtheorem{corollaries}{Corollario}[theorem]
\theoremstyle{plain}       \newtheorem{proposition}{Proposizione}[subsection]

% ---- Header e Footer ---- 
\pagestyle{fancy}
\fancyhf{}
\fancyhead[LE,RO]{\AnnoAccademico}
\fancyhead[RE,LO]{\TitoloCorso}
\fancyfoot[RE,LO]{\rightmark}
\fancyfoot[LE,RO]{\thepage}
\renewcommand{\headrulewidth}{.5pt}
\renewcommand{\footrulewidth}{.5pt}
\newcolumntype{P}[1]{>{\centering\arraybackslash}p{#1}}

% ==========================================
%   TEMPLATE PAGINE DIVISORIE (RICORRENTI)
% ==========================================
\newcommand{\PaginaDivisoria}[1]{
    \newpage
    \thispagestyle{empty}
    \vspace*{\fill}
    \begin{center}
        {\color{mainRed}\hrule height 2pt}
        \vspace{1cm}
        {\Huge \bfseries \scshape \color{mainRed} #1 \par}
        \vspace{1cm}
        {\color{mainRed}\hrule height 2pt}
    \end{center}
    \vspace*{\fill}
    \newpage
}

% Scorciatoie per le pagine che hai richiesto
\newcommand{\CoverEsercitazioni}{\PaginaDivisoria{Esercitazioni}}
\newcommand{\CoverApprofondimenti}{\PaginaDivisoria{Approfondimenti}}
\newcommand{\CoverEserciziGenerali}{\PaginaDivisoria{Esercizi Generali}}

% ==========================================
%             INIZIO DOCUMENTO
% ==========================================
\begin{document}

\begin{titlepage}
    \begin{figure}[t]
        \centering
        % Se l'immagine c'è la carica, altrimenti spazio vuoto (evita errori)
        \IfFileExists{\LogoImmagine}{\includegraphics[width=0.98\textwidth]{\LogoImmagine}}{}
    \end{figure}
    \vspace{20mm}

    \begin{Large}
     \begin{center}
        \textbf{\Dipartimento} \\
        \vspace{20mm}
        {\LARGE{\InfoCorso}}\\
        \vspace{10mm}
        
        % TITOLO ROSSO
        {\huge \bfseries \color{mainRed} \TitoloCorso}\\
        {\small{\bf \Sottotitolo}}\\
    \end{center}
    \end{Large}

    \vspace{36mm}

    % Blocchi Professori e Autori
    \begin{minipage}[t]{0.47\textwidth}
        \Professori
    \end{minipage}
    \hfill
    \begin{minipage}[t]{0.47\textwidth}\raggedleft
        \Autori
    \end{minipage}

    \vspace{25mm}

    % Linea finale ROSSA
    {\color{mainRed}\hrule height 1pt}

    \vspace{5mm}

    \centering{\large{\bf \AnnoAccademico }}

\end{titlepage} % Carica il file frontespizio.tex modificato

\tableofcontents
\newpage

% Rimuovi \maketitle standard perché usiamo frontespizio personalizzato
% \maketitle 

\begin{center}
    \vspace{-20pt}
    {\color{mainRed}\rule{11cm}{.1pt}} % Linea rossa anche qui
\end{center}

% --- I TUOI CAPITOLI ---
\section{Insiemi Numerici e Intervalli}

\subsection{Gli Insiemi Numerici Fondamentali}
Definiamo i principali insiemi numerici utilizzati in analisi matematica: % [cite: 4, 5, 7, 8]

\begin{itemize}
    \item $\mathbb{N} = \{0, 1, 2, 3, \dots\}$: Insieme dei numeri naturali (interi non negativi). % [cite: 5, 6]
    \item $\mathbb{Z} = \{ \dots, -2, -1, 0, 1, 2, \dots\}$: Insieme dei numeri interi (positivi e negativi). % [cite: 7, 22]
    \item $\mathbb{Q} = \left\{ \frac{p}{q} : p, q \in \mathbb{Z}, q \neq 0 \right\}$: Insieme dei numeri razionali (frazioni). % [cite: 8, 9]
    \item $\mathbb{R}$: Insieme dei numeri reali (include $\mathbb{Q}$ e i numeri irrazionali come $\sqrt{2}$). % [cite: 33, 35]
\end{itemize}

Vale la relazione di inclusione: 
$$ \mathbb{N} \subset \mathbb{Z} \subset \mathbb{Q} \subset \mathbb{R} $$

\subsection{Intervalli di $\mathbb{R}$}
\begin{definition}[Intervallo]
    Un sottoinsieme $I \subset \mathbb{R}$ si dice \textbf{intervallo} se, presi comunque $x, y \in I$ con $x < y$, ogni $z$ tale che $x < z < y$ appartiene a $I$. 
    In altre parole, un intervallo non ha "buchi" al suo interno. 
\end{definition}

\begin{example}
    L'insieme $A = \{x \in \mathbb{R} : 1 \le x < 2 \text{ oppure } 6 < x < 7 \}$ \textbf{non} è un intervallo. Infatti, se prendiamo $x = 1.5 \in A$ e $y = 6.5 \in A$, esiste $z=4$ tale che $x < z < y$ ma $z \notin A$. 
\end{example}

\subsubsection{Notazione degli intervalli}
Dati $a, b \in \mathbb{R}$ con $a < b$: 
\begin{itemize}
    \item \textbf{Intervallo chiuso:} $[a, b] = \{ x \in \mathbb{R} : a \le x \le b \}$ 
    \item \textbf{Intervallo aperto:} $(a, b) = \{ x \in \mathbb{R} : a < x < b \}$ (indicato talvolta come $]a, b[$) 
    \item \textbf{Intervalli semi-aperti:} 
    $$ [a, b) = \{ x \in \mathbb{R} : a \le x < b \} \quad \text{e} \quad (a, b] = \{ x \in \mathbb{R} : a < x \le b \} $$
\end{itemize}

Intervalli illimitati (semirette): 
\begin{itemize}
    \item $[a, +\infty) = \{ x \in \mathbb{R} : x \ge a \}$ (Semiretta chiusa)
    \item $(a, +\infty) = \{ x \in \mathbb{R} : x > a \}$ (Semiretta aperta)
    \item $(-\infty, b] = \{ x \in \mathbb{R} : x \le b \}$
    \item $(-\infty, +\infty) = \mathbb{R}$ 
\end{itemize}


% --- ESEMPIO PAGINE RICORRENTI ---
% Quando ti serve, scommenta queste righe:

\CoverEsercitazioni
% \input{capitoli/mie_esercitazioni}

% \CoverApprofondimenti
% \input{capitoli/miei_approfondimenti}

% \CoverEserciziGenerali
% \input{capitoli/cheatsheet}

\end{document}